
\documentclass[a4paper]{article}
\usepackage[pdftex]{graphicx}
\usepackage[utf8]{inputenc}
\usepackage{enumerate}
\usepackage{icomma}
\usepackage{amssymb}
\usepackage{polynom}
\usepackage{tikz}
\usepackage{href-ul}
\hypersetup{
	colorlinks=true,
	linkcolor=blue,
	urlcolor=blue}
\usepackage{geometry}
\geometry{a4paper, top=10mm, left=15mm, right=15mm, bottom=15mm,
headsep=0mm, footskip=12mm}


\begin{document}
\begin{tikzpicture}(10,3)
	\draw[ultra thick](2,0) --(10,0) -- (10,1.3) --(2,1.3) -- (2,0);
	\draw[fill=black](2,0)-- (0,.6) -- (2,.6) -- (2,0);
	\node at (6,.6) {\large Lösungsblatt von \href{https://www.okuyakl.de}{www.okuyakl.de}};
\end{tikzpicture}
\vspace{0.5 cm}

\noindent{\bf Aufgabe 1.}\\
a) $$(x+7)(x+11)$$
b) $$2x^2-14x \stackrel{\rm AUS}{=} 2x(x-7)$$
c) Beim Faktorisieren mit der Lösungsformel muss der Faktor vor dem $x^2$ berücksichtigt werden:
 $$x_{1/2}=\frac{2\pm \sqrt{4-2\cdot (-{3\over 2})}}{4} \Rightarrow \quad x_1 = -{1\over 2};\quad x_2={3\over 2} \Rightarrow f(x)= 2\cdot(x+{1\over 2})(x-{3\over 2})$$  
d) die Polynomdivision liefert:
$$
\polyset{style=C,div=:,vars=x}
\polylongdiv{2x^3-9x^2+7x+6}{x-3}
$$
Wir Faktorisieren mit der Lösungsformel, der Faktor vor dem $x^2$ ist wieder $2$:
$$ f(x)=(x-3)(2x^2-3x-2)\stackrel{\textnormal{LÖS}}{=}2 \cdot (x-3)(x+0,5)(x-2)$$
\noindent{\bf Aufgabe 2.}\\
a) Wir substituieren $u=x^2$ und zerlegen daraufhin:
$$f(x)=u^2-13u+36 \stackrel{\textnormal{VIE}}{=}(u-4)(u-9) \Rightarrow u_1= 4;\quad u_2=9$$
$$ \textnormal{Rücksubstitution:}\quad x_1=2\qquad x_2=-2 \qquad x_3=3 \qquad x_4=-3$$
$$f(x)=(x+2)(x-2)(x+3)(x-3)$$
b) Hier bietet sich das Ausklammern an, da überall $x$ vorkommt. Anschließend wenden wir die Lösungsformel an:
$$x^2(2x^2+x-10)=0 \Rightarrow x_{1/2}=0 $$
$$ x_{3/4}=\frac{-1\pm \sqrt{1-4\cdot 2\cdot (-10)}}{2\cdot 2} \quad \Rightarrow \quad x_3=-2,5 \quad x_4=2$$
$$f(x)=2x^2(x+2,5)(x-2)$$
c) Die Polynomdivision ergibt:
$$
\polyset{style=C,div=:,vars=x}
\polylongdiv{x^3+x^2-102x+360}{x+12}
$$
Nun kommt der Satz von Vieta zur Geltung:
$$f(x)=(x+12)(x^2-11x+30)=(x+12)(x-5)(x-6)$$
d) zuerst klammern wir $x^2$ aus und erhalten:
$$f(x)=x^2\cdot(5x^4-30x^2-275)=0 \qquad x_{1/2}=0$$
in dem zweiten Term substituieren wir $u=x^2$:
$$5u^2-30u-275=0 \qquad | :5 $$
Nun können wir diese quadratische Gleichung lösen (LÖS; VIE):
$$u^2- 6u - 55= 0 =(u-11)\cdot(u+5)$$
Rücksubstitution liefert:
$$x^2=11 \Rightarrow x_3=\sqrt{11}\, ; \quad x_4=-\sqrt{11}$$
 
$$x^2=-5 \quad \textnormal{keine weitere Lösung}$$ 
e) Zuerst klammern wir $x$ aus:
$$x\cdot(x^4+6x^3+5x^3-24x-36)=0\,;\quad x_1=0$$
Dann unterwerfen wir den Term zweier Polynomdivisionen:
$$
\polyset{style=C,div=:,vars=x}
\polylongdiv{x^4+6x^3+5x^2-24x-36}{x-2}
$$
$$
\polyset{style=C,div=:,vars=x}
\polylongdiv{x^3+8x^2+21x+18}{x+2}
$$
Der erhaltene Term erinnert an eine binomische Formel:
$$x^2+6x+9=(x+3)^2$$
Also sind die weiteren Lösungen:$$x_{4/5}=-3$$ 

\begin{center}
	\includegraphics[width=7 cm]{../../viecher/endcomic.pdf}
	
	Hier geht es zurück zum \href{https://www.okuyakl.de/math/m11nstL174/aa174.pdf}{Aufgabenblatt}
\end{center}

\end{document}